\documentclass[openany,fontset=fandol]{ctexbook}
\usepackage{amsmath,amssymb,amsthm,enumitem}
\usepackage[unicode,hidelinks]{hyperref}
\setlength{\emergencystretch}{2em}
\raggedbottom
\DeclareMathOperator{\Spec}{Spec}
\DeclareMathOperator{\Supp}{Supp}
\newcommand{\FF}{\mathcal F}
\newcommand{\GG}{\mathcal G}
\newcommand{\OO}{\mathcal O}
\newcommand{\calP}{\mathcal P}
\begin{document}
\title{代数几何基础第五卷预备笔记\\法文覆盖审计：坏集维数与截面平坦性}
\author{第98批\quad 法、英、中文三方校订}
\date{2026年11月24日}
\maketitle

\section*{范围与接续}
本批核对法文物理28页及29页开头划销段的续尾；英编稿对应物理19页末、20页上部。校订4.2、4.3的补记、命题7.2和7.3。法29从命题7.4起待续。本批未改1158页合订本。

\section*{7.1 末尾的补记：允许一部分坏点}
关于4.2和4.3，至少应在备注中给出略一般的形式。给定整数 $m$，令 $E$ 为满足下述条件的 $\xi\in P^\vee$：除去某个维数不超过 $m$ 的闭子集后，$\GG_\xi$（或 $Y_\xi$）满足 $\calP$。对于由局部良点性质定义的坏集，这意味着 $\calP$-坏点集 $Z$ 的闭包维数不超过 $m$。\footnote{法28明写“闭子集”及 $\le m$；英19末的表达易被误解为余维。对于几何整性等并非逐点定义的性质，不能无定义地使用这个等价说法。空坏集按 $\dim\varnothing=-\infty$ 处理。}

\begin{enumerate}[label=\alph*)]
\item $E$ 是 $P^\vee$ 的可构造子集。
\item 在 $S=\Spec k$ 时，若 $\FF$（或 $X$）除维数不超过 $m+1$ 的集合外满足 $\calP$，并满足所引用的泛截面定理的其余适用条件，则 $E$ 包含非空开集。
\end{enumerate}
\noindent 原稿这一段是对4.2、4.3的改写建议，并未逐项给出证明；上面的维数记法首先适用于逐点定义且坏点集可构造的性质。任意底下，可构造性按相应有限表示条件下的局部意义理解。

\section*{命题7.2：模层版本}
假设 $X$ 在 $S$ 上固有，$\FF$ 有限表示且相对 $S$ 平坦。令 $\calP$ 是4.2的(i)至(v)之一，$m$ 为整数。给定 $s\in S$，设
\[
 Z_s=\{x\in X_s:\FF_s\text{ 在 }x\text{ 处不满足 }\calP\},
 \qquad \dim\overline{Z_s}\le m+1.
\]
沿用相应泛截面结论的适用条件。\footnote{法、英7.2本身只列出性质编号，没有重述4.2(v)的特征零或非分歧限制。正文保留必要的适用范围，不把省略的限制当成已经消除；后附纯不可分覆盖反例。}

若 $k(s)$ 无限，则存在 $s$ 在 $S$ 中的开邻域 $U$ 和参数射影丛在 $U$ 上的截面 $\xi$，使得对每个 $t\in U$：
\begin{enumerate}[label=(\roman*)]
\item $Y_{\xi(t)}$ 上 $\GG_{\xi(t)}$ 不满足 $\calP$ 的点集，其闭包维数不超过 $m$；
\item 截面方程 $\phi_{\xi(t)}$ 对 $\FF_t$ 正则。
\end{enumerate}
在这些条件下，$Y_\xi$ 上的模 $\GG_\xi$ 相对 $U$ 平坦。若 $k(s)$ 未假设无限，则可先作7.1所述的有限 étale 换底，再作相应构造。\footnote{法、英底本均误写 $s$ 在 $E$ 中的邻域。$\xi$ 是基底 $S$ 的局部参数截面，故 $U\subset S$；良参数在 $P^\vee$ 中。法28首处设定点是 $s$、后处量化点也写 $s$，这里用 $t$ 分开，并统一纤维下标。}

\section*{命题7.3：概形版本及其范围}
作者称相应的概形版本本质相同：不再单列 $\FF$，改假设 $X\to S$ 平坦，并参照4.3的性质。保留 $k(s)$ 特征为零\emph{或} $f:X\to P$ 为浸入这一条件。若讨论(v)，手稿加上如下条件：对每个 $t\in S$ 及 $X_t$ 的每个不可约分支 $Z$，都有 $\dim\overline{f_t(Z)}\ge2$。\footnote{法28的性质末编号叠改，英译排作(v)，且同行手增确称“在(v)情形”；这里只明确记下这一分支的附加条件。法28左缘可辨(v)、12.2.1及若干子项，余字受黑边遮损；不复造完整旁注。}

对4.3中逐点定义且能应用第5节开性判据的(i)至(iv)，可按7.2的局部坏集方法处理：以 $\OO_X$ 代替环境模、以 $\OO_{Y_\xi}$ 代替截面余核。4.3(v)的几何整性是整体性质，不能直接套用同一个逐点坏集公式。出版的 EGA~IV 第三分册12.2.1(x)给出固有平坦有限表示族的几何整性开性；(xi)涉及无嵌入关联素循环及总重数上界。该出处说明这些整体条件另有开性定理，\emph{不}自动证明对维数容错后的7.3(v)。\footnote{正式印本12.2.1跨印179--180页；(x)、(xi)见印180。原英编注试读这两个子项，所引正式定理可核，但旁注的其余意义仍未定。}

\section*{译校补证：从单纤维选取截面}
\paragraph{可构造与开性不可混用.}
在统一关联截面 $Y\to P^\vee$ 上，若坏轨迹可构造，可在局部有限表示情形按有限个局部闭片分解；对每片应用纤维维数的可构造性，各纤维坏集的维数为各片维数的最大值。这说明维数条件怎能给出可构造集合。对于泛纤维，先核性质在原坏集以外由泛截面定理传递，并核其余假设。含泛点的参数好集可构造时才推出它含非空开集。

\paragraph{维数至多降一，并非必须恰降一.}
考虑原坏集闭包的每个不可约分支。若其像只在一个闭点上，泛超平面避开该像；否则泛截面不含整个分支，交集是其真闭子集，维数至多降一。故在继承性质的前提下，原来维数至多 $m+1$ 的坏集，在泛截面上可缩到维数至多 $m$。截面方程的正则性要另外验证，可用泛超平面避开有限多个关联点像；不从坏集维数界直接推出正则性。

\paragraph{为何可缩小 $U$ 并得到模平坦.}
在第5节的平坦有限表示假设下，先取“全纤维正则”的开参数域；在此域上，关联截面的余核相对参数域平坦。对于具备相应纤维开性定理的局部性质，截面余核的总坏轨迹是闭的。由于关联截面相对参数域固有，坏轨迹亦固有，因而其纤维维数上半连续。在 $s$ 的纤维里选到满足维数界与正则性的参数后，按7.1提升局部截面，再缩到这个开参数域。将平坦余核沿 $\xi$ 换底即得 $\GG_\xi$ 相对 $U$ 平坦。这里依赖坏轨迹闭性；不能把一般可构造集直接当作固有闭子概形。纤维维数参见 \href{https://stacks.math.columbia.edu/tag/05F6}{Stacks Project，§37.30}。

\paragraph{省去7.2(v)的限制确会出错.}
令 $k$ 为特征 $p>0$ 的代数闭域，取 $S=\Spec k$、$X=\mathbf P_k^2$、$\FF=\OO_X$，以
$[x_0:x_1:x_2]\mapsto[x_0^p:x_1^p:x_2^p]$ 为 $f$。
原概形光滑，故初始几何约化坏集为空。任一 $k$-超平面的拉回方程为
\[
 a_0x_0^p+a_1x_1^p+a_2x_2^p
   =\bigl(a_0^{1/p}x_0+a_1^{1/p}x_1+a_2^{1/p}x_2\bigr)^p .
\]
这是 $p$ 重直线，截面余核沿整条直线都不几何约化；坏集维数为1。即使取 $m=0$，维数结论也失败。方程在原概形上是正则的，因此反例不是由取消正则条件造成的。这个例子也说明单独的正则乘法、固有和平坦不足以给出(v)。

\section*{法28末至法29开头：划销的另一条证明思路}
旧稿原先指出，(v)、(vi)并非纯局部，故先把它们排除；然后建议让所得截面族平坦并有可分纤维，再用 Stein 分解中的有限 étale 因子讨论几何连通分支数的局部常值，并以此试图收回(v)。法稿的这一整段已划销；靠页边的部分字词有损，不据此补写证明。即使几何连通分支数等于一，也不能仅凭几何约化便推出几何整；例如两条相交的直线约化且连通，却非不可约。下一条命题7.4不属于本批。
\end{document}

